% report.tex — Skeleton for a scientific report on Böðn's developmental foundations
%
% Compile with: ./build.sh (or: latexmk -pdf report.tex)
%
% This is a SKELETON — sections marked [TODO] need expansion with original writing.
% Citations are real and verifiable; verify DOIs before final submission.

\documentclass[12pt, a4paper]{article}

% --- Packages ---
\usepackage[utf8]{inputenc}
\usepackage[T1]{fontenc}
\usepackage[english]{babel}
\usepackage{geometry}
\geometry{margin=2.5cm}
\usepackage{setspace}
\onehalfspacing
\usepackage{graphicx}
\usepackage{booktabs}
\usepackage{hyperref}
\hypersetup{colorlinks=true, linkcolor=blue, citecolor=blue, urlcolor=blue}
\usepackage{natbib}
\bibliographystyle{apalike}
\usepackage{enumitem}

% --- Project name ---
% \dh produces the eth character (ð) via T1 fontenc; \Bodn{} ensures
% the trailing space is not consumed after the macro.
\newcommand{\Bodn}{B\"o\dh{}n}

% --- Title ---
\title{Developmental Foundations of Tangible Interactive Learning Systems\\
       for Early Childhood:\\
       \large The \Bodn{} Project}
\author{
  Mathias Axell\\
  {\small Independent researcher}\\
  {\small \texttt{[email]}}
}
\date{\today}

% =============================================================================
\begin{document}
\maketitle

% --- Abstract ---
\begin{abstract}
\Bodn{} is an open-source, battery-powered tangible learning system designed for
children aged 3--7. The device combines physical controls (buttons, rotary
encoders, toggle switches) with multimodal feedback (LED animations, synthesized
audio, screen graphics) to deliver play experiences that target core executive
functions---working memory, inhibitory control, and cognitive flexibility---as
well as sensorimotor, cognitive, language, and social-emotional development.
This paper reviews the developmental science that informs \Bodn{}'s design,
maps each feature to established developmental domains, and identifies areas
for future development. We argue that tangible interfaces offer advantages over
touchscreen-only systems for early childhood learning, and that game-based,
guided play can meaningfully support executive function development during the
critical preschool window.

\medskip
\noindent\textbf{Keywords:} executive functions, tangible user interfaces,
early childhood, guided play, physical computing, open-source hardware
\end{abstract}

\tableofcontents
\newpage

% =============================================================================
\section{Introduction}
\label{sec:introduction}

The preschool years (ages 3--7) represent a critical period for the development
of executive functions (EF)---the cognitive control processes that enable
goal-directed behaviour \citep{diamond2013executive}. A large body of evidence
links early EF to school readiness \citep{blair2007relating}, social competence,
and lifelong well-being \citep{moffitt2011gradient}. Unlike IQ, which is
relatively stable, executive functions are highly malleable during early
childhood and respond to targeted practice
\citep{diamond2013executive, montroy2016development}.

At the same time, the proliferation of touchscreen devices has raised concerns
about passive screen time in young children \citep{radesky2015mobile}. The
World Health Organization recommends no more than one hour of sedentary screen
time per day for children aged 2--4 \citep{who2019guidelines}. This creates a
tension: digital systems can deliver adaptive, engaging practice that physical
toys cannot, yet the medium itself may undermine the developmental benefits.

Tangible user interfaces (TUI)---physical objects that bridge the digital and
physical worlds---offer a potential resolution. Since \citet{ishii1997tangible}
introduced the concept, a growing literature has explored TUI for educational
contexts, finding advantages in engagement, spatial understanding, and
embodied learning \citep{marshall2007tangible, price2003tangibles}.

This paper presents \textbf{\Bodn{}}, an open-source tangible learning system that
combines physical controls with adaptive, game-based experiences designed around
the developmental science of early childhood. We review the theoretical
foundations, describe the system design, map each feature to developmental
domains, and identify directions for future work.

% =============================================================================
\section{Background}
\label{sec:background}

% --- 2.1 Executive Functions ---
\subsection{Executive Functions in Early Childhood}
\label{sec:ef}

\citet{miyake2000unity} proposed the influential three-factor model of executive
functions: \emph{working memory} (updating), \emph{inhibitory control}
(inhibition), and \emph{cognitive flexibility} (shifting). While originally
validated in adults, subsequent work has confirmed this structure in preschoolers,
with important developmental nuances \citep{montroy2016development}.

\textbf{Working memory} is the ability to hold and manipulate information in
mind. In 4-year-olds, typical capacity is 2--4 items \citep{diamond2013executive}.
Tasks like remembering a sequence of coloured lights and reproducing it from
memory directly exercise this capacity.

\textbf{Inhibitory control} encompasses both response inhibition (suppressing a
prepotent motor response) and interference control (resisting distracting
information). The classic ``day--night'' task, where children must say ``night''
when shown a sun card, illustrates the challenge for preschoolers
\citep{diamond2013executive}.

\textbf{Cognitive flexibility} requires disengaging from one rule set and
engaging a new one. The Dimensional Change Card Sort (DCCS) task
\citep{zelazo2006dccs} is the standard assessment: children sort cards first by
colour, then by shape. Most 3-year-olds fail the switch; most 5-year-olds
succeed. The transition period (ages 3--5) represents a critical window for
practice.

% [TODO: Expand with developmental trajectory data and EF training literature]

% --- 2.2 Guided Play ---
\subsection{Guided Play and Playful Learning}
\label{sec:guided_play}

\citet{weisberg2013guided} define guided play as adult-scaffolded activities
where children explore within a structured environment designed to make
discovery productive. Guided play occupies a middle ground between free play
(child-directed, no adult structure) and direct instruction (adult-directed,
no child agency). Evidence suggests that guided play produces deeper learning
than either extreme for concept acquisition and problem-solving
\citep{hirshpasek2015putting}.

% [TODO: Expand with specific studies on guided play and EF outcomes]

% --- 2.3 Tangible User Interfaces ---
\subsection{Tangible User Interfaces for Young Children}
\label{sec:tui}

\citet{ishii1997tangible} introduced the concept of tangible user interfaces
as a way to ``give physical form to digital information.'' For young children,
who are still developing fine motor skills and abstract thinking, the physical
affordances of tangible interfaces offer several advantages over touchscreens:

\begin{itemize}[noitemsep]
  \item \textbf{Embodied interaction}: Physical manipulation engages motor
    planning, proprioception, and spatial reasoning simultaneously
    \citep{price2003tangibles}.
  \item \textbf{Reduced cognitive load}: Physical affordances (a button looks
    pressable, a knob looks turnable) reduce the interface learning burden,
    freeing cognitive resources for content \citep{marshall2007tangible}.
  \item \textbf{Multimodal feedback}: Tactile confirmation (button click,
    encoder detent) combines with visual and auditory feedback to strengthen
    sensorimotor associations.
\end{itemize}

% [TODO: Expand with TUI studies specifically targeting preschool age]

% --- 2.4 Multimodal Feedback and Sensory Integration ---
\subsection{Multimodal Feedback and Sensory Integration}
\label{sec:multimodal}

Young children benefit from redundant, multimodal feedback---receiving the same
information through multiple sensory channels simultaneously
\citep{price2003tangibles}. In \Bodn{}, every user action produces visual
(LED/screen), auditory (tone/sound), and tactile (physical click/detent)
feedback. This redundancy:

\begin{itemize}[noitemsep]
  \item Strengthens cross-modal neural connections
  \item Provides confirmation even if one channel is missed
  \item Supports different learning preferences and sensory profiles
\end{itemize}

% [TODO: Cite specific multisensory integration studies in early childhood]

% --- 2.5 Music and Auditory Development ---
\subsection{Auditory Processing and Musical Exposure}
\label{sec:auditory}

Early exposure to structured sound environments supports auditory processing
skills that underpin language acquisition and musical development
\citep{kraus2010music, trainor2005critical}. Active sound production (as
opposed to passive listening) is particularly effective.

% [TODO: Expand with findings on sound categorisation and auditory discrimination]

% --- 2.6 Bilingualism ---
\subsection{Bilingualism and Executive Function}
\label{sec:bilingualism}

A substantial literature documents enhanced executive functions in bilingual
children, attributed to the constant need to manage two active language systems
\citep{bialystok2011reshaping}. The bilingual advantage is most pronounced for
inhibitory control and cognitive flexibility---the same skills that \Bodn{}'s
game modes target.

The bilingual advantage has been contested: \citet{paap2015bilingualism} argued
that positive results reflect publication bias rather than a robust effect.
However, in the specific context of early childhood (ages 3--5), where language
systems are still being established, the evidence for enhanced inhibitory
control remains more consistent \citep{bialystok2011reshaping}. Sweden provides
a particularly favourable setting for early bilingualism: English-language media
exposure is pervasive, and Swedish preschools increasingly introduce English
informally.

\subsubsection*{Letter Case and Early Literacy}

Research on letter-name knowledge shows that preschoolers recognise uppercase
letters before lowercase ones---uppercase knowledge predicts lowercase
acquisition with a facilitation factor exceeding 16$\times$
\citep{piasta2014letter}. However, since more than 95\% of printed text uses
lowercase forms, prolonged exposure to only uppercase delays the transition to
reading real text \citep{treiman2006case}. \Bodn{} therefore uses \textbf{title
case} (first letter uppercase, remainder lowercase) for all child-facing
labels, giving the child the familiar uppercase anchor while building lowercase
familiarity.

\subsubsection*{Dual-Language Labeling}

\citet{byersheinlein2013bilingualism} demonstrated that bilingual infants
distinguish their two languages from birth using rhythmic and phonological cues,
countering the common concern that simultaneous exposure causes confusion.
\citet{tan2024translation} showed that \emph{translation equivalents}---knowing
a word in one language---significantly increases the probability of learning the
corresponding word in the other language, especially at younger ages. Classroom
research on dual-language labeling confirms that presenting objects labelled in
both languages supports vocabulary development in both
\citep{castro2011dual}. \Bodn{}'s printed NFC cards therefore present both
Swedish and English labels simultaneously (e.g.\ ``Katt / Cat''), while the
on-device UI shows one language at a time (switchable in settings) to respect
the display's space constraints.

% --- 2.7 Pretend Play ---
\subsection{Pretend Play and Social-Cognitive Development}
\label{sec:pretend_play}

\citet{lillard2013impact} reviewed the evidence linking pretend play to
cognitive development, finding associations with theory of mind, narrative
skills, and emotional understanding. While the causal direction remains debated,
the correlational evidence is strong, and pretend play is universally observed
across cultures during the preschool years.

% [TODO: Discuss the Spaceship cockpit mode as structured pretend play]

% --- 2.8 Growth Mindset and Failure Tolerance ---
\subsection{Feedback, Failure, and Growth Mindset}
\label{sec:mindset}

\citet{dweck2006mindset} distinguishes between fixed mindset (ability is innate)
and growth mindset (ability develops through effort). The type of feedback
children receive after failure shapes which mindset they develop. Process-focused
feedback (``You tried a different strategy, good!'') promotes growth mindset;
outcome-focused feedback (``Wrong!'') promotes fixed mindset.

\Bodn{}'s feedback design follows this principle: errors produce soft, warm cues
(gentle tone, dim red glow) rather than punitive signals. The intent is to
frame failure as information, not judgement.

% [TODO: Expand with specific feedback design decisions and their rationale]

% =============================================================================
\section{System Design}
\label{sec:system}

% --- 3.1 Hardware ---
\subsection{Hardware Architecture}
\label{sec:hardware}

\Bodn{} is built around an ESP32-S3 microcontroller with the following physical
interface components:

\begin{itemize}[noitemsep]
  \item 8 mini push buttons (colour-coded, via MCP23017 I2C expander)
  \item 5 arcade buttons with PWM-controlled LEDs (via PCA9685)
  \item 2 toggle switches (via MCP23017)
  \item 3 rotary encoders with push buttons (navigation, volume, throttle)
  \item 2.8'' ILI9341 TFT primary display (240$\times$320, landscape)
  \item 1.8'' ST7735 TFT secondary display (128$\times$160, portrait)
  \item 2$\times$ WS2812 8-LED sticks + 92-LED perimeter strip (RGB)
  \item MAX98357A I2S amplifier with 3W speaker
  \item INMP441 I2S MEMS microphone
  \item 6600\,mAh LiPo battery with thermal monitoring (DS18B20$\times$2)
\end{itemize}

% [TODO: Add hardware diagram or reference to docs/hardware.md]

% --- 3.2 Software Architecture ---
\subsection{Software Architecture}
\label{sec:software}

The firmware is written in MicroPython and follows a modular architecture:

\begin{itemize}[noitemsep]
  \item \textbf{Rule engines} (\texttt{*\_rules.py}): Pure logic modules with
    no hardware dependencies, testable via standard pytest on the host.
  \item \textbf{Screen modules} (\texttt{ui/*.py}): Render game state to
    displays and LEDs; handle input events.
  \item \textbf{Audio engine}: 16-voice native C mixer running on a dedicated
    CPU core, with priority channels (music, SFX pool, UI feedback) and
    built-in tone generators (sine, square, sawtooth, noise).
  \item \textbf{Session manager}: Enforces play time limits, breaks, and
    quiet hours.
  \item \textbf{i18n system}: All user-facing strings localised
    (Swedish default, English available).
\end{itemize}

% [TODO: Architectural diagram]

% =============================================================================
\section{Developmental Analysis}
\label{sec:analysis}

This section maps each \Bodn{} feature to established developmental domains and
analyses coverage. Table~\ref{tab:matrix} presents the full feature--domain
matrix.

% --- 4.1 Feature Mapping ---
\subsection{Feature--Domain Mapping}
\label{sec:mapping}

% [TODO: Expand each feature description with specific design decisions
%  linked to developmental literature]

\subsubsection{Simon (H\"armapa)}
The Simon memory game directly targets \textbf{working memory} capacity. The
starting sequence length of 2 items is calibrated to the lower bound of
4-year-old WM capacity \citep{diamond2013executive}. Each successful round
adds one item, providing adaptive difficulty within the child's zone of
proximal development \citep{vygotsky1978mind}.

\subsubsection{Rule Follow (Regelleken)}
Modelled after the DCCS task \citep{zelazo2006dccs}, Rule Follow alternates
between ``match'' and ``opposite'' rules. The ``opposite'' rule requires
\textbf{inhibitory control} (suppressing the prepotent matching response), while
the rule switch exercises \textbf{cognitive flexibility}. The frequency of rule
switches (every 4--6 correct answers) is within the range used in DCCS research.

\subsubsection{Mystery Box (Mysterium)}
An open-ended \textbf{cause-and-effect} exploration mode with no failure state.
The design follows the ``guided play'' model \citep{weisberg2013guided}: the
environment is structured (specific button--colour mappings, discoverable
combinations) but the child's exploration is entirely self-directed.

\subsubsection{Fl\"ode}
A \textbf{spatial reasoning} and \textbf{problem-solving} puzzle that requires
aligning gaps in vertical walls. The snap-to-grid mechanic reduces motor
precision demands while preserving the cognitive challenge.

\subsubsection{Garden of Life (Tr\"adg\aa rden)}
A \textbf{pattern recognition} mode based on Conway's Game of Life. The child
plants coloured seeds and observes emergent behaviour---an implicit introduction
to complex systems and self-organisation. The mode exercises no executive
functions but provides rich visual stimulation and cause-and-effect learning.

\subsubsection{Soundboard (Ljudbr\"ada)}
An \textbf{auditory processing} and \textbf{open exploration} mode.
Customisable sound banks allow parents to create themed experiences. The volume
encoder provides fine motor practice with immediate auditory feedback.

\subsubsection{Spaceship Cockpit (Rymdskepp)}
The most developmentally comprehensive mode, targeting \textbf{cognitive
flexibility} (5 scenario types requiring different inputs),
\textbf{listening comprehension} (TTS-driven instructions), and
\textbf{imaginative play} (structured pretend-play as a spaceship captain).
The adaptive difficulty system adjusts timeout windows based on the child's
recent performance, maintaining an appropriate challenge level.

\subsubsection{Story Mode (Saga)}
A data-driven branching narrative engine where the child listens to narrated
scenes and makes choices via arcade buttons to steer the plot. Primary targets:
\textbf{listening comprehension} (must listen to narration to choose),
\textbf{bilingual exposure} (full narration in both Swedish and English), and
\textbf{imaginative play} (guiding a character through a narrative world).
Stories are defined as script files on the SD card; adding new stories requires
no code changes.  The first story adapts Beatrix Potter's \emph{The Tale of
Peter Rabbit} (public domain) into a choose-your-own-adventure with four endings.

\subsubsection{High-Five Friends (Heja-kompisar)}
A reaction-time game targeting \textbf{hand-eye coordination} and
\textbf{inhibitory control}. Animals ``pop up'' on illuminated arcade buttons
and the child must tap the lit button before it disappears. The reaction window
shrinks from 2\,s to 0.5\,s as the child succeeds, providing adaptive
difficulty within the child's zone of proximal development
\citep{vygotsky1978mind}. Gentle miss feedback (soft sound, dim animation) and
automatic restart after three misses support \textbf{persistence} while keeping
emotional stakes low \citep{dweck2006mindset}. The mode also exercises
\textbf{auditory processing} (distinct hit/miss sounds) and
\textbf{self-regulation} (managing the frustration of near-misses).

\subsubsection{Sortera --- NFC Classification Game}
A tangible implementation of the Dimensional Change Card Sort (DCCS) task
\citep{zelazo2006dccs}, targeting \textbf{cognitive flexibility} and
\textbf{inhibitory control}. The device announces a sorting rule (``Find the
animals!''), and the child scans NFC cards. After several correct matches the
rule switches to a different dimension (``Find all the red ones!''), requiring
the child to inhibit the previous sorting strategy. Sortera is the first mode
to provide \textbf{active bilingual reinforcement}: on each correct scan, the
screen displays both the Swedish and English word (e.g.\ ``Katt / Cat'')
matching the physical card, and bilingual TTS announces both words---primary
language first, then secondary \citep{tan2024translation}. This dual-label
approach leverages the translation-equivalent bootstrapping effect while keeping
game instructions in the configured language for clarity.

\subsubsection{Räkna --- Progressive Math with Tangible Numbers}
A six-level NFC math mode targeting \textbf{pattern recognition} (subitising),
\textbf{sequential thinking} (counting and equation order),
\textbf{problem solving} (composing equations step by step), and
\textbf{fine motor} coordination (handling multiple cards per equation). Räkna
implements Bruner's Concrete--Pictorial--Abstract progression
\citep{bruner1966instruction,willingham2017manipulatives}: a physical dot-pattern
card is placed on the reader (concrete), the screen renders matching dots
(pictorial), and from Level~4 onwards the numeral and operator cards are
introduced (abstract). Dice-face dot patterns at low levels rather than
numerals reflect the developmental evidence that subitising is the foundation
of number sense, not numeral recognition
\citep{clementssarama2014learning}. Levels span free exploration (ages
3--4), quantity matching and more/less comparison (4--5), addition and
subtraction with tokens (5--6), and symbolic equations (6--7+). A left-to-right
linear number path on the display implements Siegler and Ramani's
\citeyearpar{siegler2009playing} finding that board-game-style number lines
measurably improve preschool number competence. Bilingual TTS announces each
quantity in Swedish and English, carrying the translation-equivalent
bootstrapping effect \citep{tan2024translation} into a cognitively engaging
math task. The mode is always self-paced and never timed --- a deliberate
design choice given the evidence that timed math exercises produce anxiety
even in preschoolers. Together with Sortera, Räkna addresses the early-math
domain that \citet{duncan2007school} identified as the single strongest
predictor of later academic achievement.

\subsubsection{Blippa --- Free-Play NFC Discovery}
A goal-less, failure-less companion to the other NFC modes targeting
\textbf{open exploration}, \textbf{cause and effect}, and \textbf{imaginative
play}. Tapping any owned card---from any card set the device knows about---
produces a full-screen emoji, a bilingual label, and either a procedural POS
``blip'' or an optional per-card sample. There is no rule to follow, no score,
and no wrong scan: the mode exists specifically to turn a growing collection
of physical cards into a low-friction exploration toy. This design is
motivated by evidence that children's causal learning is strongest when
exploration is self-directed rather than instructed \citep{bonawitz2011double},
and it intentionally creates a surface for the cashier and transit-gate
pretend play that 3--5-year-olds gravitate to, linking open exploration with
imaginative play. Unknown card identifiers within a known card set produce a
deterministic ``mystery blip'' whose pitch is derived from a stable hash of
the id, so a given card always sounds the same---reinforcing object
permanence and rehearsal of cause-effect contingencies.

\subsubsection{Sequencer (Sequenser)}
A loop-based step sequencer targeting \textbf{working memory},
\textbf{sequential thinking}, \textbf{pattern recognition}, and
\textbf{creative expression}. Five percussion tracks (arcade buttons) and one
melody track (eight mini buttons mapping to a pentatonic scale) are recorded
into an 8- or 16-step grid during looped playback. Toggle switches control
play/pause, step count, and metronome---making abstract musical concepts
tangible. The mode introduces computational thinking through creative play
\citep{papert1980mindstorms}: building patterns, running them, hearing something
unexpected, and debugging by toggling individual steps. BPM is adjustable via
encoder (70--140), providing fine motor practice with immediate auditory
feedback.

% --- 4.2 Coverage Matrix ---
\subsection{Coverage Matrix}
\label{sec:coverage}

\begin{table}[htbp]
\centering
\caption{Feature--developmental domain coverage matrix. %
  $\bullet\bullet\bullet$ = primary target;
  $\bullet\bullet$ = significant;
  $\bullet$ = incidental;
  blank = not applicable.}
\label{tab:matrix}
\small
\begin{tabular}{lccccccccccccc}
\toprule
\textbf{Domain} & \rotatebox{70}{Simon} & \rotatebox{70}{Rule Follow}
  & \rotatebox{70}{Mystery Box} & \rotatebox{70}{Fl\"ode}
  & \rotatebox{70}{Garden} & \rotatebox{70}{Soundboard}
  & \rotatebox{70}{Spaceship} & \rotatebox{70}{Story Mode}
  & \rotatebox{70}{High-Five} & \rotatebox{70}{Sequencer}
  & \rotatebox{70}{Sortera} & \rotatebox{70}{R\"akna}
  & \rotatebox{70}{Blippa} \\
\midrule
Working memory       & $\bullet\bullet\bullet$ & $\bullet\bullet$ & & $\bullet$ & & & $\bullet\bullet$ & $\bullet\bullet$ & $\bullet$ & $\bullet\bullet\bullet$ & $\bullet$ & $\bullet\bullet$ & \\
Inhibitory control   & $\bullet$ & $\bullet\bullet\bullet$ & & & & & $\bullet\bullet$ & & $\bullet\bullet$ & & $\bullet\bullet$ & $\bullet$ & \\
Cognitive flexibility & & $\bullet\bullet\bullet$ & & & & & $\bullet\bullet\bullet$ & & & $\bullet\bullet$ & $\bullet\bullet\bullet$ & $\bullet$ & \\
\midrule
Cause \& effect      & & & $\bullet\bullet\bullet$ & & $\bullet\bullet$ & $\bullet\bullet$ & & $\bullet\bullet$ & $\bullet\bullet$ & $\bullet\bullet$ & $\bullet\bullet$ & $\bullet\bullet$ & $\bullet\bullet\bullet$ \\
Pattern recognition  & $\bullet\bullet$ & & $\bullet\bullet$ & & $\bullet\bullet\bullet$ & & & & & $\bullet\bullet\bullet$ & $\bullet\bullet\bullet$ & $\bullet\bullet\bullet$ & $\bullet$ \\
Spatial reasoning    & & & & $\bullet\bullet\bullet$ & & & & & & & & $\bullet$ & \\
Sequential thinking  & $\bullet\bullet\bullet$ & & & $\bullet\bullet$ & & & & $\bullet\bullet$ & & $\bullet\bullet\bullet$ & & $\bullet\bullet\bullet$ & \\
Problem solving      & & & & $\bullet\bullet\bullet$ & & & $\bullet$ & & & $\bullet\bullet$ & & $\bullet\bullet$ & \\
\midrule
Fine motor           & $\bullet$ & $\bullet$ & $\bullet$ & $\bullet\bullet$ & $\bullet$ & $\bullet$ & $\bullet\bullet$ & $\bullet$ & $\bullet$ & $\bullet\bullet$ & $\bullet\bullet$ & $\bullet\bullet$ & $\bullet\bullet$ \\
Hand--eye coord.     & $\bullet\bullet$ & $\bullet\bullet$ & & $\bullet\bullet$ & & & $\bullet\bullet$ & & $\bullet\bullet\bullet$ & $\bullet\bullet$ & & $\bullet$ & $\bullet$ \\
Auditory processing  & $\bullet\bullet$ & & & & & $\bullet\bullet\bullet$ & $\bullet\bullet$ & $\bullet\bullet\bullet$ & $\bullet\bullet$ & $\bullet\bullet$ & $\bullet\bullet$ & $\bullet\bullet$ & $\bullet\bullet$ \\
Visual processing    & $\bullet\bullet$ & $\bullet\bullet$ & $\bullet\bullet\bullet$ & $\bullet\bullet$ & $\bullet\bullet\bullet$ & & $\bullet\bullet$ & $\bullet\bullet$ & $\bullet\bullet$ & $\bullet\bullet$ & $\bullet\bullet$ & $\bullet\bullet\bullet$ & $\bullet\bullet$ \\
\midrule
Listening comp.      & & & & & & $\bullet$ & $\bullet\bullet\bullet$ & $\bullet\bullet\bullet$ & & & $\bullet$ & $\bullet\bullet$ & $\bullet$ \\
Bilingual exposure   & $\bullet$ & $\bullet$ & & $\bullet$ & $\bullet$ & $\bullet$ & $\bullet$ & $\bullet\bullet\bullet$ & $\bullet$ & $\bullet$ & $\bullet\bullet$ & $\bullet\bullet$ & $\bullet\bullet$ \\
\midrule
Self-regulation      & $\bullet\bullet$ & $\bullet\bullet\bullet$ & & & $\bullet$ & & $\bullet\bullet$ & $\bullet$ & $\bullet\bullet$ & & $\bullet\bullet$ & $\bullet$ & \\
Persistence          & $\bullet\bullet$ & $\bullet$ & & $\bullet\bullet\bullet$ & & & $\bullet$ & $\bullet$ & $\bullet\bullet$ & $\bullet\bullet$ & $\bullet$ & $\bullet\bullet$ & \\
Imaginative play     & & & & & & & $\bullet\bullet\bullet$ & $\bullet\bullet\bullet$ & & $\bullet\bullet$ & & & $\bullet\bullet$ \\
\midrule
Open exploration     & & & $\bullet\bullet\bullet$ & & $\bullet\bullet\bullet$ & $\bullet\bullet\bullet$ & & & & & $\bullet$ & $\bullet$ & $\bullet\bullet\bullet$ \\
Divergent thinking   & & & $\bullet\bullet\bullet$ & & $\bullet\bullet$ & & & & & $\bullet\bullet\bullet$ & & $\bullet$ & $\bullet$ \\
Creative expression  & & & & & $\bullet\bullet$ & $\bullet\bullet\bullet$ & & & & $\bullet\bullet\bullet$ & & & \\
\bottomrule
\end{tabular}
\end{table}

% --- 4.3 Gap Analysis ---
\subsection{Gap Analysis}
\label{sec:gaps}

The matrix reveals strong coverage for executive functions, visual processing,
open exploration, and---with the addition of Sequencer and High-Five
Friends---auditory processing, hand-eye coordination, and persistence.
Story Mode significantly improved three previously identified gaps: listening
comprehension, bilingual exposure, and imaginative play each now have two
primary-level features (Spaceship + Story Mode). The Sequencer filled the
creative expression gap, and High-Five Friends added a dedicated hand-eye
coordination mode. The two NFC modes that followed---Sortera and R\"akna---add
a tangible manipulation layer and address two previously
underserved domains: Sortera provides a tangible Dimensional Change Card Sort
implementation \citep{zelazo2006dccs} that strengthens cognitive flexibility
and inhibitory control, while R\"akna opens the early-mathematics
domain that \citet{duncan2007school} identified as the single strongest
predictor of later academic achievement. Remaining areas with limited coverage
include:

\begin{enumerate}[noitemsep]
  \item \textbf{Spatial reasoning}: Still primarily Fl\"ode; R\"akna's linear
    number path contributes incidentally. A tangram or construction mode would
    add breadth.
  \item \textbf{Problem solving}: Fl\"ode (primary), Sequencer (debugging),
    and R\"akna (equation building) now all contribute; a scavenger-hunt style
    NFC puzzle mode would further broaden the domain.
  \item \textbf{Bilingual exposure}: Story Mode provides full narration in
    both languages, and Sortera and R\"akna add active bilingual reinforcement
    (dual-language labels and TTS on each card scan). A dedicated
    language-switching game could further leverage the bilingual advantage for
    EF development \citep{bialystok2011reshaping}.
\end{enumerate}

% =============================================================================
\section{Discussion}
\label{sec:discussion}

% [TODO: Discuss the following themes]

\subsection{Tangible Interfaces vs.\ Touchscreen Applications}
% Advantages of physical controls for this age group.
% Compare with tablet-based EF training (e.g., Cogmed, PBS games).
% Discuss the embodied cognition perspective.

\subsection{Adaptive Difficulty and the Zone of Proximal Development}
% How Böðn's difficulty scaling relates to Vygotsky's ZPD.
% The Spaceship mode's adaptive timeout system as a case study.
% Simon's sequence length progression.

\subsection{Feedback Design and Growth Mindset}
% How error feedback was designed to support a growth mindset (Dweck, 2006).
% Contrast with commercial children's games that use punitive feedback.

\subsection{Session Limits and Healthy Technology Use}
% How parental controls relate to AAP and WHO screen time guidelines.
% The distinction between ``active screen time'' and ``passive screen time.''
% Ethical considerations for children's technology design.

\subsection{Limitations}
% No formal user studies or experimental validation yet.
% Single-user focus: no cooperative play features.
% Hardware constraints: limited display resolution, no haptic feedback.
% The bilingual advantage literature is debated (cite Paap et al., 2015).

% =============================================================================
\section{Future Work}
\label{sec:future}

\begin{enumerate}
  \item \textbf{NFC tangible interaction layer}: Integrate a PN532 NFC reader to
    enable tagged physical objects (cards, stickers, figurines) as input.
    Research on tangible user interfaces consistently shows that physical
    manipulation benefits learning in preschoolers through reduced cognitive load,
    embodied engagement, and support for classification tasks
    \citep{marshall2007tangible, antle2013embodied, manches2012tangibles}. Planned
    NFC modes include:
    \begin{itemize}[noitemsep]
      \item \emph{Sortera} --- a tangible classification game implementing the DCCS
        task \citep{zelazo2006dccs} with physical cards, targeting cognitive
        flexibility and inhibitory control.
      \item \emph{Saga Builder} --- tangible storytelling where scanned
        character/setting cards trigger narration, extending the existing Story
        Mode with a richer physical input method
        \citep{sylla2012tinkrbook}.
      \item \emph{Vocabulary Explorer} --- bilingual word cards for active
        language switching, leveraging the bilingual EF advantage
        \citep{bialystok2011reshaping}.
      \item \emph{Memory Match} --- a physical search-and-select task that adds
        motor planning to working memory challenges
        \citep{manches2012tangibles}.
      \item \emph{R\"akna} --- a progressive math mode using dot-pattern NFC
        cards as tangible number tokens. Follows Bruner's
        Concrete-Pictorial-Abstract progression
        \citep{bruner1966instruction}: physical tokens (concrete)
        $\to$ dot displays (pictorial) $\to$ numerals (abstract). Six
        levels spanning subitising, counting, comparison, addition, and
        subtraction (ages 3--7). Grounded in evidence that early math
        is the strongest predictor of later academic achievement
        \citep{duncan2007school}, that linear number board games
        produce measurable and persistent gains in preschoolers'
        number competence \citep{siegler2009playing}, and that
        kindergarten number sense predicts math achievement through
        third grade \citep{jordan2009early}.
    \end{itemize}
    The NFC interaction pattern (find object $\to$ bring to reader $\to$ hold)
    gives children physical agency in a way that fixed buttons cannot, and has
    been commercially validated by products like Toniebox and Yoto Player for
    children as young as 3. See the companion document
    \texttt{nfc\_tangible\_learning.md} for the full research summary.

  \item \textbf{Empirical validation}: Conduct pre--post studies comparing
    EF measures (DCCS, digit span) in children using \Bodn{} vs.\ control
    activities.
  \item \textbf{Record \& Replay mode}: Implement voice recording and playback
    with effects (pitch shift, echo) to strengthen creative expression,
    auditory processing, and bilingual exposure.
  \item \textbf{Cooperative play}: Add turn-taking and collaborative modes to
    support social-cognitive development.
  \item \textbf{Adaptive difficulty modelling}: Implement Bayesian difficulty
    estimation to maintain each child in their zone of proximal development
    across play sessions.
  \item \textbf{Longitudinal tracking}: Use the existing session logging
    infrastructure to track developmental progress over months and years.
  \item \textbf{Active bilingual features}: Develop language-switching games
    that leverage the bilingual EF advantage.
\end{enumerate}

% =============================================================================
\section{Conclusion}
\label{sec:conclusion}

% [TODO: Summarise the argument]

\Bodn{} demonstrates that tangible, open-source hardware can deliver
developmentally informed play experiences for young children. By grounding each
feature in established developmental science---particularly the executive
function literature---and by providing physical controls rather than touchscreen
interaction, the system aims to combine the adaptability of digital systems with
the developmental benefits of hands-on play. The open-source design invites
collaboration from educators, developmental psychologists, and the maker
community to extend and validate the approach.

% =============================================================================
% References
\bibliography{report}

\end{document}
